% Job posting: https://ca.linkedin.com/jobs/view/a-i-engineer-machine-learning-engineer-at-magna-international-4461242402
% =====================================================================
%  CV - Nishal Stanislaus Silva, PhD
%  Target: Magna International / Deco Automotive (Cosma) - A.I Engineer (Machine Learning Engineer)
%  Location: Etobicoke, ON (GTA - no relocation)
%  Variant: V5 (Computer Vision / Industrial)
% =====================================================================

\documentclass[11pt]{article}
\usepackage[left=0.7in,top=0.7in,right=0.7in,bottom=0.7in]{geometry}
\usepackage{enumitem}
\usepackage[hidelinks]{hyperref}
\usepackage{titlesec}
\usepackage{array}
\usepackage{setspace}
\usepackage{needspace}
\usepackage[T1]{fontenc}
\usepackage{newtxtext,newtxmath}
\ifdefined\pdfgentounicode
  \input{glyphtounicode}
  \pdfgentounicode=1
\fi
\setstretch{1.05}
\pagenumbering{gobble}
\titleformat{\section}{\Large\scshape}{}{4em}{}[\vspace{-0.7em}\rule{\linewidth}{0.4pt}\vspace{-0.4em}]
\titlespacing*{\section}{0pt}{1em}{0.4em}
\newenvironment{cvrole}[5]{%
  \par\noindent\textbf{\large #1} | \textit{\small #2, #3}\hfill \textit{\small #4 - #5}\par
  \begin{itemize}[leftmargin=2em, itemsep=-0.3em, topsep=-0.5em]%
}{%
  \end{itemize}\par\noindent\mbox{}\par\vspace{0.1em}%
}
\newcommand{\cvName}{Nishal Stanislaus Silva}
\newcommand{\cvEmail}{nishal.silva@hotmail.com}
\newcommand{\cvPhone}{+1 613 200 2030}
\newcommand{\cvCity}{Toronto, ON}
\newcommand{\cvSite}{https://nishal.xyz}
\newcommand{\cvLinkedIn}{https://www.linkedin.com/in/nishal-silva}
\newcommand{\cvGitHub}{https://github.com/ns2max}
\newcommand{\cvStatus}{Canadian Permanent Resident, Eligible to work in Canada without sponsorship}
\newcommand{\cvTagline}{Computer Vision \& ML Engineer | Industrial Inspection, IoT \& Edge Deployment}

\begin{document}
\pagestyle{empty}
\begin{center}
    {\LARGE \scshape \textbf{\cvName}}{\large \textit{, PhD}}\\[1pt]
    {\small \cvTagline}\\[1pt]
    {\footnotesize\textit{\cvStatus}}\\[1pt]
    {\small
    \href{mailto:\cvEmail}{\cvEmail} \,\,|\,\, \cvPhone \,\,|\,\, \cvCity
    \,\,|\,\, \href{\cvSite}{nishal.xyz} \,\,|\,\, \href{\cvLinkedIn}{linkedin.com/in/nishal-silva}
    \,\,|\,\, \href{\cvGitHub}{github.com/ns2max}}
\end{center}

\section*{Professional Summary}

Computer vision and ML engineer with a PhD (University of Trento, 2025) and 10+ years across industry and academia, taking machine learning models from prototype to production on the factory floor. At MAS Holdings I spent 5.5 years building and deploying automated visual inspection and defect detection for manufacturing at scale: OpenCV pipelines feeding camera, sensor and Industrial IoT streams into live workflows, with conveyor and PLC reject integration that cut inspection time 99.5\%. I built the real-time ingestion and ETL pipelines behind that data across three continents, ran root-cause and statistical analysis on large manufacturing datasets, and introduced CI/CD, code review and pytest. My doctoral work applied the same discipline to model efficiency (14 ms inference on a Raspberry Pi 4, F1 0.76, 74$\times$ faster than baseline, JAES 2026), and Forestpin added production anomaly-scoring services for compliance and risk. Automated visual inspection is directly adjacent to weld quality inspection; the welding, stamping and assembly process specifics I would learn on the plant floor, as I learned apparel manufacturing.

\section*{Core Competencies}

\textbf{Computer Vision:} automated visual quality inspection, defect and anomaly detection, production OpenCV pipelines, template and feature matching, image registration and warping, OCR, explainable defect overlays, industrial and microscopic camera arrays, CNN and RNN architectures\\
\textbf{ML Engineering:} production-ready model development, statistical analysis and predictive modeling, predictive quality and condition/anomaly modeling, root-cause investigation, model monitoring and automated retraining, benchmark and dataset design, metrics tied to OEE, First Time Quality and scrap, TensorFlow, PyTorch, scikit-learn\\
\textbf{Systems \& Deployment:} Industry 4.0 and Industrial IoT data systems, PLC communications and conveyor integration (familiar with OPC-UA), manufacturing data pipelines, real-time ingestion and ETL, edge inference on ARM and Raspberry Pi, C++17 real-time engines, Docker, AWS, Git, pytest, CI/CD

\section*{Technical Stack}

\textbf{Languages:} Python, C++17, C, SQL, JavaScript/TypeScript, MATLAB, Shell\\
\textbf{ML \& Data:} TensorFlow, Keras, PyTorch, scikit-learn, NumPy, Pandas, SciPy; statistical modeling, anomaly detection, time-series forecasting; SQL databases; AWS (EC2, S3, ECS, RDS); Docker, Airflow, Jenkins, Git, pytest, CI/CD; dashboard and reporting tools (willing to adopt Power BI)\\
\textbf{Computer Vision:} OpenCV, template and feature matching, OCR, image registration and warping, corner detection, explainable overlays, industrial and microscopic cameras, projector-camera systems\\
\textbf{Manufacturing \& IoT:} Industrial IoT ingestion and ETL, PLC interfacing, conveyor and reject integration, smart-sensor retrofits, embedded Linux, real-time inference budgeting, familiar with OPC-UA and manufacturing data systems

\section*{Education}
\textbf{Ph.D - Information Engineering and Computer Science} \textit{\small{ - University of Trento, Italy}} \hfill \textit{Jan 2025}\\[2pt]
\noindent\textit{\small Thesis: Embedded Real-time Musical Pattern Detection for Smart Musical Instruments}\\[3pt]
\noindent\textbf{M.Sc - Telecommunication and Electronic Engineering} \textit{\small{ - Sheffield Hallam University, UK}} \hfill \textit{Aug 2018}\\[3pt]
\noindent\textbf{B.Eng (Hons) - Electronic Engineering} \textit{\small{ - Sheffield Hallam University, UK}} \hfill \textit{Mar 2014}

\section*{Professional Experience}

\begin{cvrole}{Research Engineer}{MAS Holdings (Pvt) Ltd}{Colombo, Sri Lanka}{Jan 2015}{Jul 2020}
    \item Built, trained and deployed end-to-end computer-vision and Industrial IoT systems for apparel manufacturing at scale, taking model-based inspection from prototype to production on the factory floor; operational efficiency improved by up to 300\%. Introduced CI/CD, code review and pytest; mentored engineers and ran group-wide computer-vision training.
    \item Shipped automated care-label quality inspection: OpenCV template and feature matching with an explainable defect overlay, iterated from scanner to industrial camera to a conveyor line with PLC reject; inspection time cut 99.5\%.
    \item Built loom-side fabric defect detection with a microscopic camera array (overlapping fields of view across 60-inch rolls, 5 to 6 second alert budget); pilot cut operator headcount 50\%.
    \item Built real-time ingestion and ETL pipelines for high-frequency IoT and equipment time-series across three continents; supported root-cause investigations with statistical and ML analysis of large manufacturing datasets.
    \item Delivered the Promptly on-demand garment-printing system: designed the camera-to-warp geometry-alignment algorithm, personally built the RIP integration into the automation loop, reverse-engineered PLC trigger points, and consulted on the flipping mechanism; now a Twinery/MAS commercial product across four countries.
\end{cvrole}

\needspace{5\baselineskip}
\begin{cvrole}{Postdoctoral Researcher}{University of Trento}{Trento, Italy}{Jan 2025}{Dec 2025}
    \item Owned the full ML pipeline for streaming audio and IMU/gesture data: acquisition, feature engineering, labeling, training and evaluation under domain shift, edge deployment on Raspberry Pi 4 with production monitoring and automated retraining; EU Horizon EIC Pathfinder project MUSMET.
    \item Delivered real-time inference at 14 ms on a Raspberry Pi 4, F1 0.76, 31.4\% CPU, 74$\times$ faster than the 1038 ms baseline (JAES 2026); designed the frozen-protocol benchmark suites and ablations that quantified the accuracy, latency and CPU trade-offs.
\end{cvrole}

\needspace{5\baselineskip}
\begin{cvrole}{Visiting Researcher}{McGill University (IDMIL / CIRMMT)}{Montreal, Canada}{May 2024}{Aug 2024}
    \item Fused inertial (IMU) and audio streams through a hierarchical finite state machine, reaching F1 0.78 across a 500-event corpus (IEEE IS\textsuperscript{2} 2025); profiled compute and power for real-time feasibility on self-powered embedded hardware.
\end{cvrole}

\needspace{5\baselineskip}
\begin{cvrole}{Technical Consultant}{Forestpin (Pvt) Ltd}{Colombo, Sri Lanka}{Aug 2020}{Feb 2021}
    \item Designed and deployed SQL and Python anomaly-detection scoring services for financial forensics, compliance monitoring and risk detection on large structured enterprise datasets; the same anomaly-scoring pattern transfers to predictive quality and predictive maintenance on production and equipment data.
\end{cvrole}

\section*{Selected Projects \footnotesize {\normalfont{(full portfolio: \href{https://nishal.xyz/research}{\textit{nishal.xyz/research}})}}}
\begin{itemize}[leftmargin=1.2em, itemsep=-0.25em]
    \item \textbf{Automated care-label QC (MAS):} multilingual label verification via OpenCV template and feature matching with an explainable overlay; three iterations from flatbed scanner to industrial camera to conveyor with PLC reject; 15-20 min per label down to 1-2 min per batch, +300\% throughput, fully local processing for compliance.
    \item \textbf{Loom-side fabric defect detection (MAS):} microscopic camera array with overlapping fields of view across 60-inch rolls and a 5-6 second alert budget; pilot cut operators 50\%.
\end{itemize}

\enlargethispage{4\baselineskip}
\vspace{-0.3em}
\section*{Publications \& Recognition}
\begin{itemize}[leftmargin=1.2em, itemsep=-0.25em]
    \item 10+ peer-reviewed papers (JAES, IEEE IS\textsuperscript{2}, IEEE I3DA, ACM Audio Mostly, DAFx, IEEE Asilomar), incl.\ \textit{Real-Time Audio Pattern Detection for Smart Musical Instruments}, J.\ Audio Eng.\ Soc.\ 74(3), 2026; four open-access Zenodo datasets (9,800+ recordings, 70+ musicians).
    \item COVID-19 remote-vitals computer-vision device deployed across hospital wards (GMOA recognition, Sri Lanka); MIDI Innovation Awards 2023 Finalist; Nebula, an open-source C++17 audio-feature library with ARM latency benchmarks; 2026 side projects built with AI-assisted workflows including Claude Code.
\end{itemize}

\end{document}
